\chapter{Patterns: render props and slots}
\label{ch:slots}

\begin{aside}
Some widgets need to customise their inner rendering. Render
props (a callback that returns an element) and slots (a named
inner element) are the two patterns.
\end{aside}

\section{Render prop}

A widget accepts a callable that returns the inner element:

\begin{lstlisting}
auto list = virtual_list(items,
    [](Item& i) { return text(i.name); });
\end{lstlisting}

The list owns rendering structure; the caller owns row
content.

\section{Named slots}

A widget accepts named elements:

\begin{lstlisting}
auto card = Card{}
    .header(text("Title"))
    .body(text("Body"))
    .footer(buttons);
\end{lstlisting}

Each slot is positioned by the card's layout.

\section{Default slots}

A slot may have a default if unspecified:

\begin{lstlisting}
auto card = Card{}.body(content);
// header defaults to empty, footer to empty
\end{lstlisting}

\section{Composition}

Render props compose with slots: a slot can be a render
prop that produces an element from data.

\section{When to use which}

\begin{itemize}[leftmargin=*]
  \item Many similar items (list rows): render prop.
  \item Fixed structure with varying parts (card,
    dialog): slots.
\end{itemize}

\section{Anti-patterns}

\begin{itemize}[leftmargin=*]
  \item Too many slots: just compose with column.
  \item Render prop that ignores its argument: just take
    an Element directly.
\end{itemize}

\section{Recap}

\begin{recap}
\begin{itemize}[leftmargin=*]
  \item Render prop = callback returning element.
  \item Slots = named child elements.
  \item Pick based on whether structure or content varies.
\end{itemize}
\end{recap}

\section*{Workshop}

\begin{enumerate}[leftmargin=*]
  \item Refactor a list to use a render prop.
  \item Build a card widget with header/body/footer slots.
  \item Audit your code; identify slot misuse.
\end{enumerate}
